\begin{prob}

    In lecture, we saw the following algorithm for computing a median:
    \inputminted{python}{\thisdir/include/median.py}
    We said that this algorithm has quadratic time complexity, which means that the 
    time it takes to run grows like $n^2$, where $n$ is the size of the input list.
    Let's see if this is indeed the case by timing the function on inputs of various
    size.

    \begin{subprobset}
        \begin{subprob}
            Write a function named \python{time_median} which takes in a single
            argument, \python{n}. It should create a list of \python{n} numbers
            and call \python{median} ten times on that list, timing each call.
            Your function should return the average time it took for
            \python{median} to run over all ten calls. Provide your code.

            \emph{Hint}: 
            We saw the \python{timeit} magic functions in lecture, but 
            the \python{time} function in Python's \python{time} module is more useful
            here. To time the execution of a function call, we can write:
            \begin{minted}[autogobble]{python}
                import time
                start = time.time()
                my_function_that_i_want_to_time()
                end = time.time()
                elapsed = end - start
            \end{minted}

            \begin{soln}
                 \inputminted{python}{\thisdir/include-solution/time_median.py}

            \end{soln}
        \end{subprob}

        \begin{subprob}
            Using your function \python{time_median}, time \python{median} on inputs of
            size 100, 400, 700, 1000, 1300, 1600, and 1900. Create a plot of these times,
            where the horizontal axis measures the input size, and the vertical axis measures
            the average time taken by \python{median} in seconds.

            Include your code in your solution.
            You may use any tool you like to make the plot, though we recommend using
            \python{matplotlib}.

            \begin{soln}
                Here is some code that will compute the timings.

                \inputminted{python}{\thisdir/include-solution/run_timings.py}

                This will produce a figure like that below:

                \includegraphics[scale=.6]{\thisdir/include-solution/timing.png}
            \end{soln}

        \end{subprob}
    \end{subprobset}

\end{prob}